% ADD THIS HEADER TO ALL NEW CHAPTER FILES FOR SUBFILES SUPPORT
\documentclass[../CLthesis.tex]{subfiles}
\graphicspath{{\subfix{../images/}}}
% END OF SUBFILES HEADER

\begin{document}
\chapter{Introduction}%
\label{chap:intro}

\section{Motivation}

On July 8, 2022, former Japanese Prime Minister Abe Shinz\={o} was assassinated by a gunman motivated, as investigators later established, by a personal grievance against the Unification Church (世界平和統一家庭連合), a South Korean-founded religious organisation with deep ties to Japan's ruling Liberal Democratic Party \citep{masami2024abe}. The revelation that dozens of senior politicians had cultivated relationships with a religious organisation widely described in the Japanese press as a ``harmful cult'' (有害な宗教団体) triggered one of the most intense waves of public anti-religious discourse in postwar Japanese history. On platforms including 5ch, Twitter/X, and NicoNico Dōga, denunciations of Christianity, New Age movements, and foreign religion in general proliferated rapidly; the word \textit{kyōkai} (教会, ``church'') trended alongside slurs and calls for the expulsion of ``alien beliefs'' from Japanese society. At almost exactly the same moment, thousands of kilometres to the west, Chinese social media was processing the aftermath of a years-long crackdown on unregistered Protestant congregations that had accelerated dramatically since 2018, when Beijing's revised Regulations on Religious Affairs (宗教事务条例) came into force. The closure of the Zion Church of Beijing, one of the country's largest unregistered congregations, and the December 2018 arrest of Wang Yi, pastor of Chengdu's Early Rain Covenant Church (成都秋雨圣约教会), generated millions of Weibo and Tieba posts — a significant fraction of which framed Christianity not as a persecuted faith but as a political weapon of Western imperialism, a form of ``spiritual colonisation'' (精神殖民), or an outright cult \citep{yang2011religion, ChinaAidReport2019}.

These episodes are not isolated provocations. They represent recurring flash-points in a far broader and more persistent landscape of online hostility toward Christians in East Asian digital spaces — hostility that is simultaneously rooted in specific national histories, energised by real-time political events, and systematically amplified by the architecture of social media platforms. Yet despite its scale and social salience, this form of religious hatred has attracted almost no computational research attention. It remains unmapped, unmodelled, and largely unknown to the scholarly community that studies hate speech online. This thesis is motivated by precisely this gap: the need for the first systematic, cross-lingual computational account of anti-Christian hate speech as it manifests across Chinese, Japanese, and English online platforms.

\section{Background}

\subsection{Online Hate Speech and the Religious Dimension}

The broader problem of online hate speech has attracted intense research effort over the past decade, and with good reason. Social media platforms have dramatically extended the reach and persistence of hateful content, enabling abuse that would once have been confined to marginal spaces to circulate rapidly to mass audiences \citep{castano-pulgarin_internet_2021, matamoros-fernandez_racism_2021}. Religious hate speech in particular has been recognised as a distinct and serious category of harm: it does not merely insult individuals, but systematically degrades entire communities defined by their faith, often foreshadowing or accompanying real-world violence \citep{bilewicz_hate_2020, paz_hate_2020}. International bodies from the United Nations to the Council of Europe have called for robust monitoring and counter-measures, and computational linguistics has responded with an expanding body of detection tools, annotated datasets, and analytical frameworks \citep{jahan_systematic_2023, ramos_comprehensive_2024}.

The distribution of this research effort, however, is deeply uneven. The dominant focus has been on Islamophobia and antisemitism — forms of religious hatred that are prevalent in Western contexts and for which large-scale English-language datasets are relatively accessible \citep{albadi_are_2018, vidgen_detecting_2020, Evolvi2018}. This concentration is understandable given the historical urgency of these problems, but it has produced a systematic blind spot. According to a landmark Pew Research Center report \citep{pew_rising_restrictions_2019}, Christians experienced government-sponsored or social harassment in 145 countries as of 2018 — more countries than any other religious group in the world. Christianity is, by any measure, the world's largest religion: with approximately 2.4 billion adherents representing nearly a third of the global population, it is also the most geographically dispersed, present in virtually every country on Earth \citep{jenkins_next_2011}. And yet computational models are rarely trained or evaluated to detect hostility directed against Christians, and qualitative research frameworks for understanding the specific rhetorical structures of anti-Christian discourse remain underdeveloped \citep{scheitle2020, Yancey2014}.

Part of the explanation lies in a structural assumption that has shaped hate speech research from its inception: that hate speech flows primarily from powerful majorities toward vulnerable minorities \citep{matamoros-fernandez_racism_2021}. In Western liberal democracies, where Christianity has historically occupied a position of cultural and institutional dominance, this assumption seems to place anti-Christian sentiment outside the canonical definition of marginalising hate speech and, by extension, outside the priority concerns of the field. The sociological literature has begun to complicate this picture. \citet{Yancey2014} document substantial anti-Christian bias among educated secular Americans; \citet{scheitle2020} provide survey evidence that significant proportions of American Christians report experiencing religious discrimination and hostility in everyday life. But these findings have not yet been integrated into the computational hate speech research agenda in any systematic way.

\subsection{Christianity in East Asia: A Historically Distinct Case}

The majority-assumption becomes far less tenable when the frame of reference is shifted from Western liberal democracies to East Asia. In China, Japan, and South Korea, Christianity is not a majority religion. It arrived in each country primarily through missionary activity embedded in broader colonial and geopolitical encounters: Jesuit missions in sixteenth-century Japan; Protestant missions in nineteenth-century China accompanying the Unequal Treaties; and, in Korea, missionary activity whose legacy was entangled with both the Japanese colonial period and the subsequent American military presence \citep{dunch_beyond_2002, mullins1998christianity}. In these contexts, hostility to Christianity is rarely a simple majority-backlash against a powerful institution. It is instead a historically layered phenomenon in which religious identity intersects with national sentiment, post-colonial memory, and geopolitical competition. To be a Christian in contemporary China is, in substantial portions of public discourse, to be implicitly suspect of cultural disloyalty; to be a Christian institution in contemporary Japan is to be subject to recurring accusations of foreign manipulation and cognitive coercion \citep{yang2011religion, mullins2021religion}.

This entanglement of religious and nationalist grievance has found an exceptionally fertile environment in the digital spaces of the two countries. Chinese social media platforms — including Sina Weibo, Bilibili, and Baidu Tieba — host vast communities of users who generate and circulate anti-Christian content ranging from ironic mockery of converts to coordinated harassment campaigns against visible Christians, from conspiracy theories about Vatican political influence to sustained attacks on the moral credibility of the Catholic Church in the wake of global abuse scandals. Japanese platforms — including 5ch (formerly 2channel) and NicoNico Dōga — contain long-standing threads, wikis, and comment communities dedicated to documenting and satirising what users frame as the social harms of Christian proselytising \citep{xiao_chinese_2024}. English-language platforms, meanwhile, host their own distinct tradition of anti-Christian discourse, from the sceptical and secularist communities of Reddit to more aggressively hostile communities on Twitter/X and 4chan — discourse that is culturally very different from its East Asian counterparts precisely because it takes place in a context where Christianity is still, for many participants, culturally hegemonic \citep{Yancey2014, scheitle2020}.

\subsection{Research Gap}

Despite the richness and diversity of this discourse, no computational study has systematically examined the content, targets, or rhetorical structure of anti-Christian hate speech across these three linguistic and cultural contexts. Three specific lacunae account for this gap. First, there is no publicly available annotated corpus of anti-Christian hate speech in Chinese or Japanese, and even English-language hate speech datasets rarely include religious categories beyond Islamophobia and antisemitism \citep{mathew_hatexplain_2021, mollas_ethos_2022}. Second, cross-lingual research in hate speech detection has made progress in transferring models across European language pairs, but Chinese and Japanese remain severely underrepresented in major multilingual benchmarks \citep{corazza_multilingual_2020, joshi2020state}. Third, and most fundamentally, the dominant paradigm of binary detection — classifying a post as ``hateful'' or ``not hateful'' — is structurally inadequate for the kind of analysis that would actually advance our understanding of this phenomenon. Knowing that a post is hateful tells us very little about who is being attacked, what argumentative strategy organises the attack, or what moral logic makes the attack feel justified to those who produce it \citep{elsherief_latent_2021, ollagnier_unsupervised_2023}.

This thesis addresses all three of these lacunae. By constructing new multilingual corpora, developing a cross-lingual detection pipeline, and applying a systematic unsupervised analytical framework to the resulting data, we aim to produce the first comparative, computational account of anti-Christian hate speech across Chinese, Japanese, and English online texts. The goal is not merely to identify hate speech, but to understand its fine-grained structure — to move from the question of \emph{whether} a text is hateful to the richer questions of \emph{who} is targeted, \emph{how} the targeting is accomplished, and \emph{why} the hatred takes the particular ideological forms it does.

\section{Research Questions}

The empirical and theoretical motivation developed above converges on three interlocking research questions. Together, they constitute a progressive deepening of the analysis: from the identification of targets, to the characterisation of attack strategies, to the interpretation of underlying moral logics.

The first question concerns identity. The data described above — posts on Weibo calling Christians ``spiritual colonists'', threads on 5ch cataloguing the social harms of religious organisations, English-language Reddit communities debating whether Christianity is inherently authoritarian — are not monolithic. Anti-Christian discourse attacks different entities: the institutional Church as a political actor, individual converts as traitors to national culture, clergy as moral hypocrites, believers as cognitively manipulated victims. These targets are not interchangeable. Understanding which entities bear the greatest load of hostility, and whether this distribution differs systematically across linguistic communities, is a prerequisite for any deeper analysis. We therefore ask:

\begin{enumerate}
    \item \textbf{RQ1: Who are the targets?} What entities — institutions, figures, or groups — are most frequently attacked in anti-Christian discourse across Chinese, Japanese, and English online texts, and how do these targets differ across languages?
\end{enumerate}

The second question moves from identity to rhetoric. Knowing that the Catholic Church is a frequent target in Japanese online discourse, for instance, does not tell us much about why the discourse is effective or how it is organised. Anti-Christian attacks can take the form of crude slurs and insults, but in many cases they are far more elaborate: sustained narratives alleging sexual abuse and its institutional cover-up, theories of cognitive manipulation through indoctrination and dependency, accusations of political interference and foreign allegiance, or moralistic denunciations of hypocrisy between Christian profession and practice. These rhetorical frames are not equivalent in their persuasive logic or their social function. A computational analysis that stops at identifying who is targeted has, in effect, answered only the easiest part of the question. The architecture of hate — the narrative and argumentative frames through which hostility is constructed and communicated — is both more complex and more consequential:

\begin{enumerate}
    \setcounter{enumi}{1}
    \item \textbf{RQ2: How is hatred expressed?} What are the dominant rhetorical and thematic attack frameworks (e.g., allegations of sexual abuse, cognitive manipulation, political interference, moral hypocrisy) used to articulate anti-Christian hostility, and do these patterns vary cross-linguistically?
\end{enumerate}

The third question is the most theoretically ambitious. Even after we have identified who is targeted and how the attacks are constructed, a fundamental puzzle remains: what makes these attacks feel morally legitimate — even obligatory — to those who produce them? This question becomes especially provocative when we compare the three linguistic communities in our study. In English-language contexts, where Christianity remains a culturally dominant institution with substantial political representation, anti-Christian hatred must justify itself against the grain of mainstream cultural norms; it tends to draw on a vocabulary of individual rights, rational autonomy, and resistance to institutional authority. In Chinese and Japanese contexts, by contrast, anti-Christian sentiment can often draw on the considerably more powerful rhetoric of national identity and civilisational integrity — Christianity as foreign imposition, believers as culturally disloyal, the Church as an arm of Western hegemony. Are these simply different surface expressions of the same underlying moral concerns, or do they reflect genuinely distinct moral logics? Moral Foundations Theory \citep{haidt2012righteous, graham_liberals_2009} offers a framework for operationalising this question: if anti-Christian discourse in different languages mobilises different moral foundations — harm/care versus loyalty/betrayal versus sanctity/degradation — this should be detectable through the semantic content of the texts themselves:

\begin{enumerate}
    \setcounter{enumi}{2}
    \item \textbf{RQ3: Why do people hate?} What underlying moral motivations, as characterised by Moral Foundations Theory, drive anti-Christian rhetoric in each language community, and can vector-space projection methods reveal cross-linguistic differences in these motivations?
\end{enumerate}

Taken together, these three questions constitute a \emph{who}--\emph{how}--\emph{why} framework for the analysis of hate speech that is applicable beyond the specific case of anti-Christian discourse. Rather than treating hate speech detection as a binary classification problem, we treat it as the first step in a richer programme of hate speech interpretation — one that aims to understand the structure of online hostility well enough to explain it, and ultimately to counter it more effectively.

\section{Contributions}

The main contributions of this thesis are threefold. First, it provides the first systematically constructed multilingual corpus of anti-Christian hate speech in Chinese, Japanese, and English, together with detection models fine-tuned for each language. These resources fill a significant empirical gap in the field and are designed to support future research on non-Western religious hate speech. Second, it proposes and validates a cross-lingual analytical pipeline that moves beyond binary classification to support fine-grained target identification and rhetorical frame analysis, using unsupervised methods based on multilingual sentence embeddings and BERTopic \citep{grootendorst_bertopic_2022}. Third, it demonstrates, for the first time, how Moral Foundations Theory can be operationalised at scale through vector-space projection to produce cross-linguistically comparable measurements of the ideological motivations underlying hate speech — a methodological contribution with broad applicability to the computational analysis of political and identity-based discourse.

\section{Thesis Structure}

The remainder of this thesis is organised as follows. Chapter~\ref{chap:related-work} reviews the existing literature across four dimensions: the sociology of religious hate speech and the specific case of anti-Christian hostility; the evolution of computational hate speech detection from rule-based systems to large language models; the state of cross-lingual and low-resource research with particular attention to Chinese and Japanese; and the development of unsupervised and clustering-based approaches to hate speech analysis. Chapter~\ref{chap:material} describes the data collection and annotation pipeline, including the crawling and filtering of Chinese and Japanese social media data, the construction of a multilingual religious hate speech dataset with cross-lingual annotation consistency, and the fine-tuning of transformer-based language models for hate speech detection across the three languages. Chapter~\ref{chap:methods} presents the full unsupervised analytical pipeline: the embedding of sentences into a shared multilingual vector space, the application of BERTopic for topic modelling and cluster interpretation, and the FrameAxis-based projection method used to measure moral foundations across languages. Chapter~\ref{chap:bertopic} outlines the topic structures, whereas Chapters~\ref{chap:who}--\ref{chap:why} provide a comprehensive interpretation of the results for each research question. Through these chapters, I offer cross-linguistic comparisons of attack targets, rhetorical frameworks, and moral motivations, discussing the broader implications of convergence and divergence across the three communities. Chapter~\ref{chap:conclusion} concludes with a synthesis of the main findings, a critical assessment of the limitations of the present approach, and directions for future work.

\subfilebibliography
\end{document}
